% Generated draft: requires mathematical and source review.
\documentclass[11pt]{article}
\usepackage[margin=25mm]{geometry}
\usepackage{fontspec}
\setmainfont{Latin Modern Roman}
\usepackage{amsmath,amssymb,mathtools,unicode-math}
\setmathfont{Latin Modern Math}
\usepackage{enumitem,needspace,xurl,hyperref}
\setlength{\parindent}{0pt}
\setlength{\parskip}{5pt}
\setlength{\emergencystretch}{3em}
\newcommand{\N}{\mathbb{N}}
\newcommand{\Z}{\mathbb{Z}}
\newcommand{\Q}{\mathbb{Q}}
\newcommand{\R}{\mathbb{R}}
\newcommand{\C}{\mathbb{C}}
\newcommand{\F}{\mathbb{F}}
\newcommand{\A}{\mathbb{A}}
\newcommand{\PP}{\mathbb{P}}
\newcommand{\E}{\mathbb{E}}
\newcommand{\eps}{\varepsilon}
\newcommand{\dd}{\,\mathrm{d}}
\begin{document}
\section*{UnsolvedMath ID: 30003855}
\textbf{Source problem code:} catalogue entry
\subsection{Definitions and mathematical statement}
For each fixed integer \(n\geq 3\), let \(\mathcal D_n\) denote the class of nonempty smooth bounded domains \(\Omega\subset\R^n\). Here ``domain'' is interpreted in its usual sense of a connected open set, and smoothness means \(C^\infty\) boundary. The source does not explicitly discuss disconnected open sets or specify a boundary regularity class beyond ``smooth.'' No convexity, connectedness of the boundary, or fixed topological type is required by the supplied statement.

For \(\Omega\in\mathcal D_n\), write
\[
|\partial\Omega|=\mathcal H^{n-1}(\partial\Omega)
\]
for its boundary surface measure. The first positive Steklov eigenvalue \(\sigma_1(\Omega)\) is the least positive number \(\sigma\) for which there exists a nonzero function \(u\) satisfying
\[
\begin{cases}
\Delta u=0 & \text{in }\Omega,\\
\partial_\nu u=\sigma u & \text{on }\partial\Omega,
\end{cases}
\]
where \(\nu\) is the outward unit normal. Thus the zero eigenvalue corresponding to constant functions is excluded. Define
\[
F_n(\Omega)=\sigma_1(\Omega)|\partial\Omega|^{1/(n-1)},
\qquad
S_n=\sup_{\Omega\in\mathcal D_n}F_n(\Omega),
\]
allowing \(S_n=+\infty\) unless finiteness is established.

\textbf{Problem.} For every fixed integer \(n\geq3\), determine \(S_n\) and characterize the domains \(\Omega_*\in\mathcal D_n\), if any, satisfying
\[
F_n(\Omega_*)=S_n.
\]
If the supremum is not attained, describe sequences \((\Omega_j)_{j\geq1}\subset\mathcal D_n\) such that
\[
F_n(\Omega_j)\longrightarrow S_n,
\]
with divergence to \(+\infty\) understood if \(S_n=+\infty\), and describe their limiting geometry insofar as a precise limiting formulation can be supplied.

The source asks about a domain or a degenerating sequence realizing the supremum, but specifies neither a notion of geometric convergence nor a normalization for discussing limits. Accordingly, the displayed convergence defines a maximizing sequence only; it does not assert convergence of the domains or prescribe a degeneration. The supplied open status and literature assessment are dated curation claims, not independently verified mathematical conclusions.
\subsection{Short English statement}
For each dimension \(n\geq3\), determine the largest possible value of the first positive Steklov eigenvalue multiplied by boundary area to the power \(1/(n-1)\), over all smooth bounded domains in \(\R^n\). Identify any maximizing domains or describe maximizing sequences and their possible degeneration. The source leaves the meaning of geometric degeneration unspecified; connectedness is interpreted through the usual meaning of ``domain.''
\subsection{Sources}
\begin{itemize}
\item Local UnsolvedMath registry, \texttt{data/unsolvedmath/problems.json}, supplied record with permanent UnsolvedMath ID \(30003855\), title ``Sharp Normalized Steklov Supremum in Higher Dimensions,'' exact source problem code \texttt{catalogue entry}, in the collection ``Oberwolfach Reports Open Problems.'' No individual proposer is supplied; the proposed year is recorded as 2018.
\item Supplied report attribution: \emph{Geometrie (2019)}, Oberwolfach Reports, pp.\ 1635--1684, DOI \texttt{10.4171/owr/2018/27}: \url{https://doi.org/10.4171/owr/2018/27}. The background describes Oberwolfach Reports as documenting mathematical workshops and their open-problem discussions. The supplied title year 2019 and the 2018 report identifiers are retained without reconciliation.
\item Bucur--Ferone--Nitsch--Trombetti, ``Weinstock inequality for convex domains,'' as labeled in the supplied literature sources: \url{https://arxiv.org/abs/1710.04587}.
\item Fraser--Schoen, ``Shape optimization for the Steklov problem,'' as labeled in the supplied literature sources: \url{https://www.sciencedirect.com/science/article/pii/S0001870819301410}.
\item ``Oberwolfach Report 27/2018,'' as labeled in the supplied literature sources: \url{https://ems.press/content/serial-article-files/46749?nt=1}.
\item Literature assessment checked at \textbf{2026-08-21}, according to the supplied record: the problem is labeled open; sharp Weinstock-type bounds are reported for restricted classes such as convex domains, while the unrestricted higher-dimensional sharp supremum question is reported to remain open. This is supplied dated curation derived from the OWR record, not an independently verified assessment. The supplied background expressly requests independent verification of time-sensitive status and cited claims.
\item The registry labels the statement \texttt{corrected\_verified}, with the supplied explanation that mathematical content was preserved while notation, typography, and sentence structure were repaired for publication. This label is recorded as source metadata and does not establish mathematical correctness or independent source verification.
\end{itemize}
\end{document}
